
\documentclass[12pt,a4paper]{article}

% Packages
\usepackage[utf8]{inputenc}
\usepackage{geometry}
\usepackage{amsmath}
\usepackage{graphicx}
\usepackage{booktabs}
\usepackage{hyperref}
\usepackage{xcolor}
\usepackage{fancyhdr}

% Page settings
\geometry{
    margin=1in
}

% Header and footer
\pagestyle{fancy}
\fancyhf{}
\lhead{Technical Writing}
\rhead{LaTeX Documentation}
\cfoot{\thepage}

% Hyperlink settings
\hypersetup{
    colorlinks=true,
    linkcolor=blue,
    urlcolor=blue
}

% Document begins
\begin{document}

% Title Page
\begin{titlepage}
    \centering

    \vspace*{2cm}

    {\Huge\bfseries Technical Documentation\\[0.3cm]
    Using LaTeX\par}

    \vspace{1.5cm}

    {\Large A Practical Learning Report\par}

    \vspace{2cm}

    \rule{\textwidth}{0.5pt}

    \vspace{1cm}

    {\large Technical Writing Course\par}

    \vspace{0.5cm}

    {\large Academic Session 2026--27\par}

    \vspace{1cm}

    {\large Prepared by: [Your Name]\par}

    \vfill

    {\large November 2026\par}

    \rule{\textwidth}{0.5pt}

\end{titlepage}

% Abstract
\begin{abstract}
LaTeX is a powerful document preparation system widely
used for creating academic papers, technical reports,
research articles, and mathematical documents.

This report explores the fundamentals of LaTeX, its
document structure, commonly used commands, mathematical
expressions, tables, figures, and referencing techniques.
It also highlights the importance of LaTeX in producing
professional and well-organized technical documentation.
\end{abstract}

\newpage

% Table of Contents
\tableofcontents
\newpage

% Section 1
\section{Introduction to LaTeX}

LaTeX is a document preparation system built on top of
TeX. It allows users to create professionally formatted
documents using markup commands.

Unlike traditional word processors, LaTeX separates
the content of a document from its formatting.

It is commonly used in:
\begin{itemize}
    \item Academic research papers
    \item Technical reports
    \item Books and articles
    \item Mathematical and scientific documents
\end{itemize}

\section{Basic LaTeX Commands}

LaTeX uses commands beginning with a backslash.

Some commonly used commands are:

\begin{itemize}
    \item \textbackslash section\{\}: Creates a section.
    \item \textbackslash subsection\{\}: Creates a subsection.
    \item \textbackslash textbf\{\}: Makes text bold.
    \item \textbackslash textit\{\}: Makes text italic.
    \item \textbackslash emph\{\}: Emphasizes text.
    \item \textbackslash begin\{\} and \textbackslash end\{\}:
    Define environments.
\end{itemize}

For example:

\textbf{This is bold text.}

\textit{This is italic text.}

\section{Document Structure}

A basic LaTeX document contains three major parts:

\begin{enumerate}
    \item Document class declaration.
    \item Preamble containing packages and settings.
    \item Document body containing the actual content.
\end{enumerate}

The commands \texttt{\textbackslash begin\{document\}}
and \texttt{\textbackslash end\{document\}} mark the
beginning and end of the document body.

\section{Mathematical Equations}

One of the major advantages of LaTeX is its ability
to display mathematical expressions clearly.

For example, the quadratic formula is:

\begin{equation}
x = \frac{-b \pm \sqrt{b^2 - 4ac}}{2a}
\end{equation}

The summation notation can be represented as:

\begin{equation}
S = \sum_{i=1}^{n} i
\end{equation}

These features make LaTeX especially useful for
scientific and mathematical documentation.

\section{Tables and Figures}

Tables help organize information in a structured way.

\begin{table}[h]
\centering
\caption{Common LaTeX Commands}
\begin{tabular}{ll}
\toprule
\textbf{Command} & \textbf{Purpose} \\
\midrule
\textbackslash section & Creates a section \\
\textbackslash title & Defines document title \\
\textbackslash author & Defines author name \\
\textbackslash caption & Adds a caption \\
\textbackslash label & Creates a reference label \\
\bottomrule
\end{tabular}
\end{table}

Figures can also be inserted into LaTeX documents
using the graphicx package.

\begin{figure}[h]
\centering
\fbox{
\begin{minipage}{0.7\textwidth}
\centering
\vspace{0.7cm}
\textbf{LaTeX Document Workflow}

\vspace{0.4cm}
Write Source Code

$\downarrow$

Compile with LaTeX

$\downarrow$

Generate PDF Document
\vspace{0.7cm}
\end{minipage}
}
\caption{Basic LaTeX document preparation workflow}
\label{fig:workflow}
\end{figure}

As shown in Figure~\ref{fig:workflow}, LaTeX converts
source code into a formatted document.

\section{References and Bibliography}

References are important in technical writing because
they acknowledge the sources of information used.

A bibliography can be created using the
\texttt{thebibliography} environment.

\section{Advantages and Applications}

Some important advantages of LaTeX include:

\begin{itemize}
    \item Professional document formatting.
    \item Excellent mathematical typesetting.
    \item Automatic numbering of sections and figures.
    \item Easy management of references.
    \item Consistent formatting throughout documents.
\end{itemize}

LaTeX is widely used in engineering, computer science,
mathematics, and scientific research.

\section{Conclusion}

LaTeX is an effective tool for producing structured,
accurate, and professional technical documents.

Through this practical exercise, we explored its
document structure, formatting commands, mathematical
equations, tables, figures, and referencing methods.

Learning LaTeX provides an important foundation for
creating high-quality academic and technical documents.

\begin{thebibliography}{9}

\bibitem{latex}
Leslie Lamport,
\textit{LaTeX: A Document Preparation System},
Addison-Wesley.

\bibitem{overleaf}
Overleaf Documentation,
\url{https://www.overleaf.com/learn}

\bibitem{latexproject}
The LaTeX Project,
\url{https://www.latex-project.org/}

\end{thebibliography}

\end{document}
